\subsection{紧空间的一致结构}

定义 1 拓扑空间 X 上的一致结构称为与 X 的拓扑相容，如果后者与该一致结构诱导的拓扑重合。

若拓扑空间上存在与它的拓扑相容的一致结构，则称该拓扑空间为可一致化的，并称其拓扑为可一致化的。

存在不可一致化的拓扑空间，例如任何不满足公理 \((\mathrm{O}_{\mathrm{III}})\) 的空间（§ 1 第 2 目命题 2 的推论 3）；于是就产生了在什么条件下一个拓扑空间可一致化的问题。

我们要到第九章 § 1 才给出这个问题的完整回答。在本节中我们只考察一个重要的特殊情形，即 X 是紧空间的情形。这时我们有下面的定理：

定理 1 在紧空间 X 上，恰有一个与 X 的拓扑相容的一致结构；这个一致结构的一致邻域是对角线 \(\Delta\)（在 \(\mathbf{X} \times \mathbf{X}\) 中的）的所有邻域。此外，赋有这个一致结构的 X 是完备一致空间。

定理的最后一部分是直接的；因为 X 上的每个 Cauchy 滤子都有聚点{[}公理 (C){]}，因而收敛（§ 3 第 2 目命题 5 推论 2）。

我们接着证明，若 X 上存在与 X 的拓扑相容的一致结构，则它的一致邻域组成的集 U 是 \(\Delta\) 的所有邻域组成的集。我们已知每个一致邻域都是 \(\Delta\) 的邻域（§ 1 第 2 目命题 2），因此只需证明反之每个 \(\Delta\) 的邻域都属于 U。假设存在 \(\Delta\) 的一个不属于 U 的邻域 U；则当 V 取遍 U 时，诸集合 \(V \cap CU\) 构成紧空间 \(X \times X\) 上一个滤子 G 的基；于是 G 有一个聚点 \((a, b)\)（不属于 \(\Delta\)）。由于 U 是比 G 粗的滤子，\((a, b)\) 也是 U 的聚点。但按假设由 U 定义的一致结构是 Hausdorff 的；因此 U 的诸集合的闭包之交是 \(\Delta\)（§ 1 第 2 目命题 2 的推论 2 及命题 3）；于是我们得到矛盾。

因此剩下要证明的是，\(\mathfrak{B}\)（即 \(\Delta\) 在 \(\mathbf{X} \times \mathbf{X}\) 中的邻域组成的集）是某个与 \(\mathbf{X}\) 的拓扑相容的一致结构的一致邻域集。为此只需证明 \(\mathfrak{B}\) 是 \(\mathbf{X}\) 上某个 Hausdorff 一致结构的一致邻域集；因为这时该一致结构诱导的拓扑将比 \(\mathbf{X}\) 的拓扑粗（第一章 § 2 第 2 目命题 3），从而必与后者重合（第一章 § 9 第 4 目定理 2 推论 3）。

\(\mathfrak{B}\) 显然满足公理 \((\mathbf{F}_{\mathrm{I}})\) 与 \((\mathbf{F}_{\mathrm{II}})\)。我们来证明公理 \((\mathbf{U}_{\mathrm{II}})\) 与 \((\mathbf{U}_{\mathrm{III}})\) 也得到满足，并且 \(\Delta\) 是 \(\mathfrak{B}\) 中诸集合之交。先证最后一点：每个由单独一点 \((x, y) \in \mathbf{X} \times \mathbf{X}\) 组成的集合都是闭的，因为 \(\mathbf{X}\) 是 Hausdorff 的；因此若 \(x \neq y\)，则 \((x, y)\) 在 \(\mathbf{X} \times \mathbf{X}\) 中的余集是 \(\Delta\) 的邻域。由于对称映射 \((x, y) \to (y, x)\) 是 \(\mathbf{X} \times \mathbf{X}\) 到自身的同胚，可知 \(V \in B\) 蕴含 \(\overline{V} \in B\)，由此得 \((U_{II})\)。最后，假设 \((U_{III})\) 不满足；则存在集合 \(V \in B\)，使得对所有 \(W \in B\)，集合 \(\overline{W} \cap C V\) 非空；于是当 W 取遍 B 时，诸集合 \(\overline{W} \cap C V\) 构成 \(X \times X\) 上的一个滤子基，且该滤子基有一个聚点 \((x, y)\)（不在 \(\Delta\) 中）。现在，由于 X 是正则的（第一章 §9 第 2 目命题 1 的推论），存在不相交的开邻域 \(U_{1}\)（x 的）与 \(U_{2}\)（y 的），并且存在闭邻域 \(V_{1} \subset U_{1}\) 与 \(V_{2} \subset U_{2}\)（分别为 x 与 y 的闭邻域）。令 \(U_{3} = C(V_{1} \cup V_{2})\)，并考虑邻域

\[
\mathbf {W} = \bigcup_ {i = 1, 2, 3} (\mathrm{U} _ {i} \times \mathrm{U} _ {i}) \text {of} \Delta \text {in} \mathbf {X} \times \mathbf {X}.
\]

由这些定义立即得出，若 \((u, v) \in W\) 且 \(u \in V_1\)（相应地 \(u \in U_1\)），则必有 \(v \in U_1\)（相应地 \(v \in U_1 \cup U_3 = \complement V_2\)）；因此邻域 \(V_1 \times V_2\)（\((x, y)\) 在 \(X \times X\) 中的）不与 \(W\) 相交，这就得到矛盾。证明完毕。

注 1 对每个有限开覆盖 \(\Re = (\mathrm{U}_i)_{1\leqslant i\leqslant n}\)（\(\mathbf{X}\) 的），集合

\[
\mathrm{V} _ {\Re} = \bigcup_ {i = 1} ^ {n} \left(\mathrm{U} _ {i} \times \mathrm{U} _ {i}\right)
\]

是 \(\Delta\)（在 \(\mathbf{X} \times \mathbf{X}\) 中的）的邻域，并且这些集合 \(\mathrm{V}_{\Re}\) 构成 \(\Delta\) 的一个基本邻域系（因而构成 \(\mathbf{X}\) 上唯一一致结构的一致邻域基本系）。设 \(\mathbf{W}\) 是 \(\Delta\) 在 \(\mathbf{X} \times \mathbf{X}\) 中的任一邻域；则对每个 \(x \in \mathbf{X}\)，存在开邻域 \(\mathbf{U}_x\)（\(x\) 在 \(\mathbf{X}\) 中的），使 \(\mathbf{U}_x \times \mathbf{U}_x \subset \mathbf{W}\)。由于诸 \(\mathbf{U}_x\)（ \(x \in \mathbf{X}\) ）构成 \(\mathbf{X}\) 的开覆盖，故存在有限多个点 \(x_i\)（ \(1 \leqslant i \leqslant n\) ）使诸 \(\mathbf{U}_{x_i}\)（ \(1 \leqslant i \leqslant n\) ）构成一个覆盖 \(\Re\)（\(\mathbf{X}\) 的）。于是有 \(\mathrm{V}_{\Re} \subset \mathrm{W}\)，这就证明了所述断言。

由于这个原因，X 上唯一的一致结构常称为有限开覆盖的一致结构（参看第九章 § 4 习题 17）。

推论 1 紧空间的每个子空间都是可一致化的。

推论 2 每个局部紧空间都是可一致化的。

因为由 Alexandroff 定理（第一章 § 9 第 8 目定理 4），局部紧空间同胚于紧空间的子空间。

注 2 可能发生这样的情形：与局部紧空间的拓扑相容的一致结构不止一个且互不相同。

例如，我们已经看到，在无限离散空间上，有多于一个相容于离散拓扑的不同一致结构（§ 2 第 2 目的注）。

尽管如此，与可一致化空间的拓扑相容的一致结构的唯一性并不是刻画紧空间的性质；存在不紧但其一致结构唯一的局部紧空间（习题 4）。

定理 2 每个连续映射 \(f\)（从紧空间 \(X\) 到一致空间 \(X'\) 中的）都是一致连续的。

令 \(g = f \times f\)，则 \(g\) 在 \(X \times X\) 上连续（第一章 § 4 第 1 目命题 1 推论 1）；因此对每个开一致邻域 \(V'\)（\(X'\) 的），\(\overline{g}(V')\) 是 \(X \times X\) 的开子集，它显然包含对角线。于是定理由定理 1 得出，因为 \(X'\) 的开一致邻域构成一个一致邻域基本系（§ 1 第 2 目命题 2 推论 2）。

在定理 2 的假设下，\(f\) 在任一子空间 \(A\)（\(X\) 的）上的限制是一致连续的；因此（§ 3 第 6 目定理 2）：

推论 设 A 是紧空间 X 的稠密子空间，f 是 A 到完备 Hausdorff 一致空间 X' 中的映射。则 f 可由连续性延拓到整个 X 上，当且仅当 f 一致连续。

